\documentclass[11pt]{article}
\usepackage[margin=25mm]{geometry}
\usepackage{fontspec}
\setmainfont{Latin Modern Roman}
\usepackage{amsmath,amssymb,mathtools}
\usepackage{unicode-math}
\setmathfont{Latin Modern Math}
\usepackage{xurl,hyperref}
\hypersetup{colorlinks=true,urlcolor=blue}
\setcounter{secnumdepth}{0}
\setlength{\parindent}{0pt}
\setlength{\parskip}{5pt}
\setlength{\emergencystretch}{3em}
\newcommand{\sref}[2]{\hyperlink{source:#1:#2}{[S#2]}}
\newcommand{\cataloguescope}[1]{\paragraph{Recorded target.}#1}
\newcommand{\R}{\mathbb R}
\newcommand{\C}{\mathbb C}
\newcommand{\N}{\mathbb N}
\newcommand{\Z}{\mathbb Z}
\newcommand{\Q}{\mathbb Q}
\newcommand{\F}{\mathbb F}
\newcommand{\eps}{\varepsilon}
\begin{document}
\section{Classification of Quantizations of Symplectic and Poisson Manifolds}\label{statement}

\subsection{Definitions and mathematical statement}

\paragraph{Shared notation.}$\mathbb N=\{1,2,\ldots\}$; the standard number fields have their usual meanings. All additional conventions are stated locally.


\paragraph{Definitions and conventions.}
A symplectic manifold carries a closed nondegenerate two-form; a Poisson manifold carries a bracket on functions satisfying the Leibniz and Jacobi identities. A deformation quantization is an associative product $f\star g=fg+\sum_{j\geq1}\hbar^j C_j(f,g)$ with $(f\star g-g\star f)/(i\hbar)$ equal to the Poisson bracket to leading order. Classification is relative to an equivalence notion, such as formal changes of variables preserving the classical limit.

\paragraph{Target retained from the canonical source catalogue.}
\paragraph{Catalogue wording retained for traceability.}
Given a symplectic or Poisson manifold, how can all possible quantizations be classified, including deformation quantizations of its algebra of functions?

\paragraph{Source scope and qualification.}
Formal, strict analytic, geometric and representation-theoretic quantizations are distinct. The broad question does not prescribe one equivalence relation for every such notion, and known formal classifications must not be mistaken for resolution of all requested quantizations.
\subsection{Short English statement}
Classify quantizations of symplectic and Poisson manifolds, including their deformation products.
\subsection{Research attempt}
\paragraph{Model, date and investigated scope.}
OpenAI GPT-6 Astra (configured model \texttt{gpt-6-astra}), ultra reasoning; generated 11 October 2026. Constructs a constant-Poisson formal product and proves an explicit family of equivalent first-order coefficients; does not classify arbitrary quantizations or establish analytic realizations.

\paragraph{Sources and literature check.}
The complete supplied definition, including source qualifications and the dated literature assessment, was reviewed. The calculation below is an elementary restricted attempt motivated by [S2]; it does not claim novelty or a fresh literature-status determination.

\paragraph{Context and investigated scope.} The complete supplied definition, including source qualifications and the dated literature assessment, was reviewed. The calculation below is an elementary restricted attempt motivated by [S2]; it does not claim novelty or a fresh literature-status determination.

\paragraph{Constant-Poisson construction.}
On the $d$-torus let $\Pi$ be a fixed real antisymmetric $d\times d$ matrix and define
$\{f,g\}=\sum_{a,b}\Pi_{ab}(\partial_af)(\partial_bg)$.
For Fourier modes $e_k(x)=e^{ik\cdot x}$, $k\in\mathbb Z^d$, set formally
\[
 e_k\star e_l=\exp\left(-\frac{i\hbar}{2}k^T\Pi l\right)e_{k+l}.
\]
We work with formal power series in $\hbar$ whose coefficients are trigonometric polynomials, so every coefficient of a product is a finite sum. Bilinearity gives
\[
 k^T\Pi l+(k+l)^T\Pi m
 =l^T\Pi m+k^T\Pi(l+m),
\]
which proves associativity on modes and hence on this algebra. The constant mode is a unit. Also
\[
 \frac{e_k\star e_l-e_l\star e_k}{i\hbar}
 =-k^T\Pi l\,e_{k+l}+O(\hbar^2)
 =\{e_k,e_l\}+O(\hbar^2).
\]
This verifies the required classical bracket directly, including degenerate $\Pi$.

\paragraph{Testing what is actually classified.}
Let $D$ be a differential operator preserving trigonometric polynomials with $D1=0$, and let $T=I+\hbar D$. Its inverse exists formally as $\sum_{j\geq0}(-\hbar D)^j$. Transport the product by
$f\star' g=T^{-1}((Tf)\star(Tg))$.
Associativity and the unit are preserved. If $C_1$ is the first coefficient of $\star$, expansion gives
\[
 C'_1(f,g)=C_1(f,g)+(Df)g+f(Dg)-D(fg).
\]
The added expression is symmetric in $f,g$, so the antisymmetric first-order bracket cannot distinguish these products. For example, $D=a\partial_1^2$ gives the change $-2a(\partial_1f)(\partial_1g)$ for any real $a$. Infinitely many displayed coefficient formulas can thus represent the same formal equivalence class.

\paragraph{First obstruction to a classification.}
The construction supplies one family on a constant-Poisson torus and exhibits explicit redundancy; it does not show every deformation is equivalent to one of these products. Nor does a formal inverse or formal exponential establish convergence at nonzero $\hbar$, a norm completion, positivity of states, or a classification of representations. The broader source question distinguishes formal, strict analytic, geometric and representation-theoretic quantizations. A classification statement needs a chosen category and equivalence relation before an invariant can be proved complete. No assertion that the established formal theory in [S3] is unresolved is intended.

\paragraph{Verification and adversarial audit.}
Checked associativity coefficientwise, the commutator sign against the Poisson bracket, formal inversion, the unit condition D1=0, and the symmetric gauge coboundary. Model self-audit only; no independent human proof checking. Prompts A and B were applied to this canonical record and its qualified source scope. Checked associativity coefficientwise, the commutator sign against the Poisson bracket, formal inversion, the unit condition D1=0, and the symmetric gauge coboundary. Retained unresolved status and zero original-target progress; Show completeness of an invariant for a specified quantization category and equivalence relation; this calculation neither exhausts formal products nor controls strict analytic ones.

Coordinator checked the constant-Poisson Fourier cocycle, classical commutator sign and first-order product change under formal differential equivalence. This additional check was model-assisted, not independent human verification.

\paragraph{Tools, process and contributions.}
Python 3 was used for file preparation and structural checks; no numerical experiment or theorem prover was used. Contribution prompt A was followed by prompt B. GPT-6 Astra with ultra reasoning drafted the research and performed its stated self-audit; any additional coordinator contribution is identified in the verification and audit. A GPT-6-based Codex coordinator assembled the manuscript and checked repository integration; its exact variant and reasoning configuration were not exposed. No human mathematical contribution, independent human proof check or formal proof certification is supplied. Complete derivations and the consulted source references are included here; private reasoning transcripts are not published.

\paragraph{Final status and first remaining gap.}
\textbf{UNRESOLVED}. Show completeness of an invariant for a specified quantization category and equivalence relation; this calculation neither exhausts formal products nor controls strict analytic ones.

\section{Completion Estimate}
\noindent\textbf{COMPLETION: 0\%}
The calculation is an auxiliary or restricted attempt and establishes no newly verified portion of the full source target. Show completeness of an invariant for a specified quantization category and equivalence relation; this calculation neither exhausts formal products nor controls strict analytic ones.

\subsection{Sources}
\begin{itemize}
\item[\hypertarget{source:statement:1}{S1}] ULAM.AI and the UnsolvedMath Contributors, UnsolvedMath, catalogue entry. Catalogue curation under CC BY 4.0; definitions and English exposition added locally.
\url{https://huggingface.co/datasets/ulamai/UnsolvedMath}.
\item[\hypertarget{source:statement:2}{S2}] Topologie (2016), Oberwolfach Reports, DOI 10.4171/owr/2016/35, pp. 2009--2068. Primary report accessed 2026-10-11; relevant PDF pages 15.
\url{https://doi.org/10.4171/owr/2016/35}.
\item[\hypertarget{source:statement:3}{S3}] Kontsevich, deformation quantization (catalogue-listed context; not independently consulted in this preparation).
\url{https://arxiv.org/abs/q-alg/9709040}.
\end{itemize}
\label{end:statement}
\end{document}
